%/*% pdflatex -file-line-error -interaction nonstopmode %

\documentclass{article}
\usepackage{graphicx}
\usepackage{amsmath}
\usepackage{amsfonts}
\usepackage{bookmark}
\usepackage{booktabs}
\usepackage{exercises}
\graphicspath{ {./images} }
\usepackage{indentfirst}
\usepackage{amsthm}
\usepackage{yhmath}
\usepackage{placeins}
\newcommand*{\oldneg}{\mathord{\sim}}
\newcommand\showdiv[1]{\overline{\smash{)}#1}}
\newcommand\Mydiv[1]{%
\smash{\raise.3ex\hbox{$\big)$}}\mkern-8mu
        \overline{\enspace\strut#1}}
\setlength{\tabcolsep}{5pt}
\theoremstyle{definition}
\newtheorem{definition}{Definition}

\setlength{\parindent}{.5in}
\setlength{\parskip}{12pt}
\renewcommand{\arraystretch}{1}

\title{2.7 Answers}
\begin{document}
\maketitle

\section*{1}


For the equation $(x-1)^2(x+2)=0$, a repeated solution is $1$ because the factor $x-1$ appears twice.

\section*{2}


The difference is that complex conjugates relate to complex numbers, while irrational conjugates relate to irrational numbers. In particulr, the complex 
conjugate of a complex number $a+bi$ is defined to be $a-bi$, while the irrational conjugate for a number $a+\sqrt{b}$ where $\sqrt{b}$ is irrational 
is defined to be $a-\sqrt{b}$.

\section*{6}


By the fundemental theorem of algebra, $f(z)=-7z^4+z^2-25$ has $4$ solutions.

\section*{7}


By the fundemental theorem of algebra, $g(s)=12s^7-9s^6+4s^5-s^3-20s+50$ has $7$ solutions.

\section*{8}


By the fundemental theorem of algebra, $h(x)=-x^12+7x^8+5x^4-8x+6$ has $12$ solutions. 

\section*{9}


By the fundemental theorem of algebra, the answer is $C$.


\section*{18}


Using the rational roots theorem, we can see that $-\frac{1}{2}$ is one solution. We can use synthetic division to find the other factor, which is 
\[ 4x^3+2x^2-12x-6 \] Using the rational roots theorem, we can see that $-\frac{1}{2}$ is also a solution of the above factor. We can again use synthetic 
division to find the other factor, which is \[ 4x^2-12 \] It has solutions of:
\[ 4x^2-12=0 \]
\[ x^2=3 \]
\[ x=\pm\sqrt{3} \]
Thus, we have found four zeros of the initial equation: \[ -\frac{1}{2},-\frac{1}{2},-\sqrt{3},\sqrt{3} \] Because we are dealing with an equation of degree 
four, we know that there are exactly four zeros. Thus, we can conclude that we have found all the zeros of the function.

\section*{19}


We are given the function \[ h(x)=2x^4+13x^3+19x^2-10x-24 \] and are asked to find all of it's zeros. Using the rational roots theorem, we can tell that 
$1$ is a solution, meaning we can use synthetic division to find the other factor, which is \[ 2x^3+15x^2+34x+24 \] We can use the rational roots 
theorem again to tell that it has a zero of $-2$, meaning we can again use synthetic division to find another factor, which is 
\[ 2x^2+11x+12=(2x+3)(x+4) \] Meaning the other zeros are $x=-\frac{3}{2}$ and $-4$. We have thus far found four zeros: \[ 1,-2,-\frac{3}{2},-4 \]
By the fundemental theorem of algebra, the original equation has exactly four zeros, meaning we have found all of the zeros of this equation. 
	
